\chapter{Asthma}

\section*{Definition and Overview}
Asthma is a chronic inflammatory condition of the airways that causes them to become narrow, swollen, and hypersensitive, with increased mucus production in response to triggers. This leads to recurring episodes of wheezing, coughing, chest tightness, and breathlessness. The airway narrowing in asthma is characteristically variable and reversible, either spontaneously or with treatment. Asthma cannot be cured, but in the great majority of people it can be very well controlled with appropriate medicines, correct inhaler technique, and an individualised self-management plan, allowing a normal, active life.

\section*{Epidemiology}
Asthma is one of the most common chronic diseases in the world, affecting roughly one in ten children and around one in twelve adults. It is more common in boys before puberty and in women after puberty. In Australia, asthma is particularly prevalent, partly because of high rates of atopy and environmental factors. Higher-risk groups include people with a family history of asthma or allergic disease, smokers, people who are obese, and those exposed to household or workplace triggers such as house dust mites, moulds, and industrial chemicals. Despite treatment advances, asthma continues to cause significant school and work absence and remains a cause of preventable death.

\section*{Aetiology and Pathophysiology}
\begin{itemize}
    \item \textbf{Genetic predisposition:} asthma and allergic conditions such as eczema and hay fever cluster in families
    \item \textbf{Allergens:} house dust mites, pollens, moulds, and animal dander can trigger symptoms in sensitised people
    \item \textbf{Infections:} viral respiratory infections, particularly in early childhood, are a common trigger of attacks
    \item \textbf{Irritants:} tobacco smoke, air pollution, cold air, strong smells, and certain workplace chemicals
    \item \textbf{Exercise:} exertion, especially in cold, dry air, can provoke bronchospasm
    \item \textbf{Pathophysiology:} airway inflammation and oedema, bronchial smooth muscle contraction, and excess mucus combine to narrow the airways and obstruct airflow
    \item \textbf{Remodelling:} long-standing, poorly controlled inflammation can produce permanent structural changes in the airway wall
\end{itemize}

\section*{Clinical Features}
\begin{itemize}
    \item Wheezing, a high-pitched whistling sound heard mainly when breathing out
    \item Cough, which may be dry and is often worse at night or in the early morning
    \item Chest tightness or a feeling of pressure within the chest
    \item Shortness of breath, worse on exertion or at night
    \item Symptoms that come and go, varying from day to day and week to week
    \item Trigger-related symptoms, such as after exercise, cold air, allergen exposure, or a cold
    \item During an attack: rapid breathing, difficulty speaking, and visible use of the accessory muscles of breathing
\end{itemize}

\section*{Differential Diagnosis}
Conditions that can mimic asthma include:
\begin{itemize}
    \item Chronic obstructive pulmonary disease (COPD), particularly in older smokers
    \item Vocal cord dysfunction, which causes similar breathlessness and wheeze
    \item Heart failure causing pulmonary congestion and oedema
    \item Bronchitis, bronchiectasis, and other airway infections
    \item A foreign body in the airway, especially in young children
    \item Gastro-oesophageal reflux, which can cause cough and chest symptoms
    \item Anaphylaxis and other allergic reactions
\end{itemize}

\section*{When to Seek Medical Care}
See a doctor if you or your child have a cough, wheeze, or breathlessness that keeps recurring, wakes you at night, or limits activity. A review is also needed when symptoms are not well controlled, when the reliever is being used more than two to three times per week, or after any asthma attack requiring urgent care.

\textbf{Call 000 (or local emergency services) for an asthma emergency:}
\begin{itemize}
    \item Severe breathlessness, or an inability to speak in full sentences
    \item Reliever medicine not helping or not lasting
    \item Lips or face turning blue
    \item Exhaustion, confusion, or drowsiness
    \item Gasping for breath or an inability to breathe
\end{itemize}

\section*{Diagnosis and Investigations}
\begin{itemize}
    \item \textbf{History:} the pattern of symptoms, triggers, night waking, reliever use, and any personal or family history of allergic disease
    \item \textbf{Spirometry:} the key diagnostic test, measuring the forced expiratory volume in one second; significant improvement after a reliever supports asthma
    \item \textbf{Peak flow monitoring:} a small handheld device measures airflow, and significant day-to-day variability suggests asthma
    \item \textbf{Bronchial challenge testing:} performed in a lung function laboratory when the diagnosis remains uncertain
    \item \textbf{FeNO testing:} measurement of exhaled nitric oxide as a marker of eosinophilic airway inflammation
    \item \textbf{Allergy testing:} skin-prick tests or specific IgE blood tests for common allergens
    \item \textbf{Chest X-ray:} sometimes used to exclude other lung conditions or complications
\end{itemize}

\section*{Management}
\begin{itemize}
    \item \textbf{Reliever inhalers:} short-acting bronchodilators taken as needed to open the airways quickly
    \item \textbf{Preventer inhalers:} inhaled corticosteroids taken daily to reduce airway inflammation and prevent attacks
    \item \textbf{Combination inhalers:} a preventer combined with a long-acting reliever for people who need additional treatment
    \item \textbf{Written asthma action plan:} a personalised plan describing daily preventer use and the steps to follow when symptoms worsen
    \item \textbf{Trigger avoidance:} reducing exposure to smoke, house dust mites, moulds, pollens, and other triggers
    \item \textbf{Vaccinations:} annual influenza and pneumococcal vaccination to reduce infection-related flares
    \item \textbf{Smoking cessation:} essential, together with avoidance of second-hand smoke
    \item \textbf{Severe asthma care:} specialist review, add-on medicines, and biologic therapies for selected patients with difficult asthma
\end{itemize}

\section*{Complications}
\begin{itemize}
    \item Asthma attacks requiring emergency care or hospital admission
    \item Persistent poor control with night waking and time off school or work
    \item Reduced lung function and airway remodelling from long-term undertreatment
    \item Side effects of medicines, particularly oral corticosteroids with repeated or prolonged courses
    \item Anxiety, sleep disturbance, and reduced quality of life
    \item Rarely, life-threatening attacks and death, mostly where asthma is undertreated or action plans are not followed
\end{itemize}

\section*{Prognosis and Outcomes}
Most people with asthma live active, full lives with good control. Symptoms fluctuate over time and may change with age; many children improve markedly or appear to outgrow their symptoms as they grow older. Well-controlled asthma causes few daytime symptoms and little restriction of activity. Control depends on the consistent use of preventer treatment, correct inhaler technique, and adherence to a written action plan. People with severe or difficult-to-control asthma benefit from specialist review and targeted therapies.

\section*{Prevention and Self-Care}
\begin{itemize}
    \item Use preventer inhalers daily as prescribed, even when you feel well
    \item Learn and practise correct inhaler technique with your doctor or pharmacist
    \item Follow your written asthma action plan at all times
    \item Avoid smoking and second-hand smoke, and seek help to quit
    \item Reduce exposure to known triggers such as house dust mites, pollens, and moulds
    \item Keep up to date with influenza and pneumococcal vaccination
    \item Stay active with regular exercise, using your reliever before exercise if advised
    \item Have regular asthma reviews and update your action plan
\end{itemize}

\section*{Sources and Further Reading}
The following organisations provide reliable, evidence-based information on asthma for patients and healthcare students.

\begin{itemize}
    \item NHS. \textit{Asthma: symptoms, diagnosis, and treatment.} https://www.nhs.uk/conditions/asthma/
    \item Global Initiative for Asthma (GINA). \textit{Global strategy for asthma management and prevention.} https://ginasthma.org/
    \item National Asthma Council Australia. \textit{Australian Asthma Handbook.} https://www.asthmahandbook.org.au/
    \item Asthma + Lung UK. \textit{Asthma information and support.} https://www.asthmaandlung.org.uk/
    \item MedlinePlus. \textit{Asthma.} https://medlineplus.gov/asthma.html
    \item World Health Organization. \textit{Asthma fact sheet.} https://www.who.int/news-room/fact-sheets/detail/asthma
\end{itemize}
